\section{Discussion and Conclusion}
\label{sec:assessment}

This section assesses the guarantees the design provides, the costs it incurs, and the situations in which it is most appropriate. It concludes with the broader direction the architecture is intended to support.

\subsection{Bounded Guarantees}

Each assurance the \glslink{synchronic-web}{Synchronic Web} provides comes from a specific mechanism and holds within a specific boundary. A relying \glslink{journal}{journal} can verify that retained material matches a commitment it has accepted, because \glslink{digest}{digests} over the \glslink{node}{node} structure bind every available region to that commitment (\cref{sec:integrity-state}). The claim holds only under a correct verifier, and recovering omitted material depends on custodians that still retain it. A journal can also establish that one state existed before another, since signed succession orders a journal's own \glslink{head}{heads} and inclusion orders a remote head before the local head that contains it (\cref{sec:history-entanglement}). That ordering depends on synchronization, on key continuity, and on witnesses to expose forks. When a result follows from a stored definition, any party holding the \glslink{self-encoding}{self-encoding} \glslink{object}{object} can recompute it (\cref{sec:self-encoding}), provided it runs a compatible evaluator under confinement and host limits. Finally, a destination can authenticate a request as coming from a particular journal through the exact committed \glslink{route}{route} and the interface key it reaches (\cref{sec:journal-web}), although the result itself rests on the reachable endpoint, secure transport, and application handling of replay.

Other properties lie outside the cryptographic core: domain evidence supports truth, and preservation sustains availability. Selective disclosure, destination policy, secure transport, and protected storage sustain confidentiality. Confidentiality extends into persistence, since journals hold private data and administrative credentials alongside proofs and history. Everything that holds or executes that material, from the client and evaluator to persisted state and the key store, therefore belongs to the trusted computing base. Backups in particular are live credential material. The limits of each mechanism also clarify what failure looks like. A surviving digest identifies lost evidence, two signatures can expose equivocation, and a deterministic transition can reproduce a supplied timestamp without establishing its accuracy.

\subsection{Costs and Architectural Fit}

Each property carries a corresponding cost. Integrity requires hashing, persisted handles, proof paths, and verification work, and history adds retained states and key stewardship. Partial knowledge saves content storage, yet each retained proof still carries ancestor evidence, \glslink{stump}{stumps}, and serialization metadata. Self-encoding objects bring the familiar burdens of executable code: evaluation, confinement, debugging, compatibility, and migration. Federation adds exact-object routing, direct transport, authorization local to the destination, origin verification of historical results, and application handling of replayed requests (\cref{sec:keys-principals}). The magnitude of these costs depends on the workload and deployment. Because these costs are real, architectural fit matters, as summarized in \cref{tab:architectural-fit}. The Synchronic Web becomes compelling when independent roots, selective retention, durable interpretation, historical correction, and cross-institutional evidence coincide, as they may in production traceability, portable identity, decentralized archives, or agent reputation. Applications that need a single shared decision, as in payments, settlement, or exclusive allocation, can still add coordination above the journal layer. Journals give these domain protocols retained inputs and histories, and each application defines its shared decision procedure.

\begin{table}[htbp]
  \centering
  \small
  \begin{tabularx}{\textwidth}{@{}lX@{}}
\toprule
\textbf{System} & \textbf{Best Suited For} \\
\midrule
Ordinary database & One operator querying current data efficiently. \\
Content-addressed store & Anyone fetching unchanging data by its hash. \\
Transparency log & One authority publicly accountable for its actions. \\
Conflict-free replicated data type & Replicas that must converge on the same state. \\
Global ledger & Participants agreeing on one shared transaction order. \\
Synchronic Web & Independent operators sharing distributed computational state. \\
\bottomrule
\end{tabularx}

  \caption{Architectural Fit}
  \label{tab:architectural-fit}
\end{table}

\subsection{From Architecture to Specification}

The implemented system serves as a reference for the mechanisms described here, and future work centers on a specification of their observable invariants. That specification should define them without freezing the current object layouts or service decomposition, and conformance tests should connect each invariant to accepted and rejected traces. Its assurance boundary should be equally explicit, stating that catastrophic key loss remains a governance matter (\cref{sec:control-evidence}) and that request isolation bounds a single operation rather than all workloads (\cref{sec:shared-computation}). These limits identify where applications, operators, and future specifications must add policy without weakening the integrity claims made here.

\subsection{Toward a Plural Web}

The architecture begins from a modest observation: durable information needs both continuity and a point of view. Continuity comes from structure, in the form of integrity-protected state, signed succession, retained definitions, and cross-history receipts. The point of view comes from a journal, which chooses its roots and applies its own policy to the evidence it retains and the relationships it follows. Joining the two is the deeper value of timeline entanglement in a \glslink{resource}{resource} architecture~\cite{Maniatis2002timeline}: remote history becomes locally retained evidence, self-encoding objects give that evidence meaning, and journal-relative \glslink{path}{paths} make it findable and citable. The composition requires neither a common platform nor a single authoritative state, since histories intersect while each keeps its own sequence and observers hold different portions of one committed state.

A plural Web will still contain disagreement, loss, compromise, and institutional change. The Synchronic Web does not resolve these issues, but it provides the means to preserve and connect diverse viewpoints within one network of independently held histories. As more independent operators join the network of journals, the Web can become not only a place to publish information but a structure on which to build the plural memory of human cooperation.
